\cotefile{125}
\scanpage{2}
Dualité locale et dualité \deleted{\unsure{affine}} projective
% archive sheet number: 1
\scanpage{3}
\[
\begin{aligned}
P^n&\supset X\supset Y\\
r&>n>0
\end{aligned}
\qquad H^p(X^r-Y^n,F)=0\quad\text{pour}\quad
\unsure{-p|p}\geq\unsure{r+n}
\]
\[
H^p_Y(X,F)=0\qquad\text{pour}\quad p>\qquad
\Big|\quad{}^{\displaystyle H}
\]

\uline{$X$ de dim $r$}
\diagram{\begin{tikzcd}[ampersand replacement=\&]
H^p(X^r,F)\arrow[r]\&H^p(X^r-Y,F)\arrow[r]\&H^{p+1}_Y(X^r,F)
\end{tikzcd}}
\[
\ast
\]

\usetikzlibrary{shapes.geometric}
\diagram{\resizebox{\linewidth}{!}{\begin{tikzcd}[ampersand replacement=\&,column sep=1.3em,row sep=2.7em]
|[draw]| {H^p(\deleted{\unsure{X}}X^r-Y,F)=0\quad\text{pour}\quad p\geq m}
  \arrow[r,Leftrightarrow]\arrow[d,Leftrightarrow]
 \& |[draw]| {\begin{gathered}H^p(X-Y,\underline O(\unsure{-n|+n|n}))=0\\
       \text{pour }n\text{ grand}\\p\geq m\end{gathered}}\\
|[draw]| {\begin{gathered}H^p_Y(X^r,F)\longrightarrow H^p(X^r,F).\quad
       \text{bijectif pour }p>m\\[\jot]
       [\,\text{surjectif pour }p=m\,]\end{gathered}}
  \arrow[r,Leftrightarrow]\arrow[d,Leftrightarrow,"{\text{dualité}}"{draw,ellipse}]
 \& |[draw]| {\begin{gathered}\text{Kif kif pour}\\F=\underline O(-n),\\n\text{ grand}\end{gathered}}\\
|[draw]| {\begin{gathered}
 \operatorname{Ext}^{q}_{\underline O_{\widehat X^{r}}}
  (\widehat X^r,\widehat F,\widehat\Omega^r_{X^r/k})
  \xleftarrow{\sim}
 \operatorname{Ext}^{q}_{\underline O_{X^r}}(X^r,F,\Omega^r_{X^r/k})
       \quad\text{bijectif pour }q<r-m\\
 {}[\,\text{injectif pour }\deleted{\unsure{q}}q=r-m\,]\end{gathered}}
 \arrow[dr,Leftrightarrow]\&\\
|[draw,ellipse]| {\begin{gathered}
 H^0(X^r,\underline O(n))\overset{\unsure{\text{iso.}}}{\xrightarrow{\sim}}
 H^0(\widehat X^r,\widehat{\underline O}(n))\\ \text{et}\\
 H^i(\widehat X^r,\widehat O(n))=0\quad\text{pour}\\
 n\text{ grand},\quad 0<i<r-m
 \end{gathered}}\arrow[r,equal]
 \& |[draw]| {\begin{gathered}
 H^q(X^r,\deleted{\unsure{\Omega}}{\underline O}(n))
 \longrightarrow H^q(\widehat X^r,\widehat{\underline O}(n))\\
 \left.\begin{gathered}\text{bijectif pour }q<r-m\\
 {}[\,\text{injectif pour }q=r-m\,]\end{gathered}\right\}n\ \text{grd}
 \end{gathered}}\arrow[d,Leftrightarrow]\\
 \& |[draw]| {\begin{gathered}
 H^q(X^r,F)\longrightarrow H^q(\widehat X^r,\widehat F)\quad\text{est}\\
 \left.\begin{gathered}\text{bij. pour }q<r-m\\
 {}[\,\text{inj. pour }q=r-m\,]\end{gathered}\right\vert\
 \text{pour }F\text{ loc libre}\end{gathered}}
\end{tikzcd}}}
\margin{(N.B. Le cas $p=m$ peut être omis, car $H^m(X^r,\underline O(-n))=0$ pour $n$ grand)}
% slanted
\margin{\uline{N.B.} $H^p(X^r-Y^n,\underline O(\unsure{\ast|n}))=0$ si $0<p<r-n-1$ [et on a un \deleted{\unsure{isomorphisme}} \inserted{\unsure{iso.}} $H^0(X^r,\underline O(\ast))\longrightarrow H^0(X^r-Y^{\unsure{n}},\underline O(\ast))$ si $n<r-1$]}

Cas \unsure{important} $\Rightarrow\quad [\unsure{n|N|m}=\dim Y]\qquad
 m=r-n,\quad r-m=n$
% archive sheet number: 2
\scanpage{4}
\mspage{- III-18 -}
$C_y$ ; il est immédiat que tout élément de cette algèbre peut s'écrir\unsure{e} \margin{\handwritten{$\ast$}}$d\handwritten{\otimes}(1/t)$, où $t\handwritten{\in}C$ \deleted{\illegible{1}} est tel que $t\handwritten{\notin}j_y\inserted{\handwritten{=0}}$, et $d\handwritten{\in}D$ ; elle, a par suite une base de la forme $(d_i\handwritten{\otimes}1)_{1\handwritten{\leq}i\handwritten{\leq}n}$, où $d_i\handwritten{\in}D$. Nous allons voir que l'on peut supposer\deleted{\illegible{1}} que les $d_i$, qui sont des sections de $\underline O_X$ au-dessus de $f^{-1}(U)$, sont des restrictions de sections de $\underline O_X$ \uline{au-dessus de} $X$. Remarquons pour cela que le préschéma $f^{-1}(y)$ a un espace de base formé d'un nombre fini de points $x_j$ ($1\handwritten{\leq}j\handwritten{\leq}m$ ; II,2,7,pro).17) formant un sous-espace \uline{discret} de $f^{-1}(U)$ ; l'algèbre \deleted{\illegible{1}} $D\handwritten{\otimes}_C C_y$, qui est l'anneau de ce schéma affine, est donc composée directe des sous-algèbres $\underline O_{x_j}$ (I,\uline{2,7,prop.19 et I,4,1,prop.4}), \margin{\handwritten{!}}et on peut \uline{évidemment} supposer que chacun des éléments $d_i\handwritten{\otimes}1$ est un élément de base
% slanted
\margin{\handwritten{? \unsure{contenu} dans un \unsure{des}}}de l'un des $\underline O_{x_j}$. Comme $X$ est affine, il y a pour chaque $j$ un élément $f_j\handwritten{\in\Gamma}(X,\underline O_X)=B$ tel que $x_j\handwritten{\in}X_{f_j}\handwritten{\subset}f^{-1}(U)$ \deleted{\unsure{et} $x_k\handwritten{\unsure{\notin}}X_{f_j}$ \unsure{pour} $k\handwritten{\unsure{\ne}}j$} ; en outre, on sait qu'il y a un entier $q$ tel que chacune des sections $d_i f_j^q$ (où $d_i\handwritten{\otimes}1$ est dans $\underline O_{x_j}$) se prolonge en une section $b_i\handwritten{\in}B$ (I,1,4,th.2). Comme $f_j$ est inversible au point $x_j$ et que les $d_i\handwritten{\otimes}1$ correspondant à $x_j$ sont nulles aux $x_k\handwritten{\ne}x_j$, les éléments $b_i\handwritten{\otimes}1$ forment encore une base de $D\handwritten{\otimes}_C C_y$ sur $C_y$, ce qui établit notre assertion.

Cela étant, \uline{considérons l'homomorphisme} $\handwritten{\Gamma}(Y,\underline O_Y^n)\handwritten{\to}\handwritten{\Gamma}(X,\underline O_X)=B$ \uline{faisant correspondre}
% slanted
\margin{\handwritten{\unsure{ou} \uline{\unsure{bien}} \uline{\unsure{faisant}}}}$b_i$ au i-ème élément de base $(0,..,1,..,\deleted{\unsure{0}}0)$ de $\handwritten{\Gamma}(Y,\underline O_Y^n)$ ; cet homomorphisme définit évidemment un homomorphisme de faisceaux $\boxed{u:\underline A^n\handwritten{\to}\underline B.}$ \deleted{\illegible{1}} Les faisceaux $\underline A^n$ et $\underline B$ sont quasi-cohérents sur $Y$ (II,2,2,prop.5), donc cohérents\inserted{\handwritten{ puisque $Y$ est \unsure{nœthérien|noethérien} et $B$ de type fini par hypothèse}}, et il en est dnnc de même du noyau et du conoyau de $u$. Or, par construction, si \inserted{$s$ est} une section de $\operatorname{Ker}u$ ou de $\operatorname{Coker}u$ au-dessus d'un voisinage de $y$, le germe de $s$ au point $y$ est nul, donc $s$ est nulle dans un voi\deleted{\unsure{s}}sinage \unsure{de} $y$ ; on en conclut qu'il existe un voisinage ouvert $V$ de $y$ tel que la restriction $u\handwritten{\vert}V$ soit un \uline{isomorphisme} $\underline A^n\handwritten{\vert}V\handwritten{\to}\underline B\handwritten{\vert}V$.

Cela étant, en vertu du th.2, il suffit de prouver que pour tou\handwritten{t}
% archive sheet number: 3
\scanpage{5}
\[
\deleted{X\unsure{=|\subseteq|\subset}}P^r\supset X\supset Y
\]
\[
H^p_Y(X,F)\simeq H^p_Y(P,F)=\varinjlim\operatorname{Ext}^p(P;\underline O_{Y_n},\unsure{\mathcal F|F})
\]
\[
[\,H^p(X-Y,F)=\varinjlim\operatorname{Ext}^p(P;{\unsure{J|\mathcal J}}^n,F)\,)
\]
Son dual est
\[
\begin{aligned}
&\varprojlim\operatorname{Ext}^{r-p}(P;F,\underline O_{Y_n}\otimes\Omega^r_{P^r/\unsure{k|h}})\\
&\xrightarrow{\sim}\operatorname{Ext}^{r-p}(\widehat Y,\widehat F,\widehat{\Omega^r}_{P^r})
\end{aligned}
\]

\diagram{\resizebox{\linewidth}{!}{\begin{tikzcd}[ampersand replacement=\&,column sep=.65em]
H^p_Y(X,F)\&{\unsure{\mathord{-}|\to}}\&H^p(X,F)\arrow[r,dash]\&
 H^p(X-Y,F)\arrow[r,dash]\&H^{\unsure{p-1|p+1}}_Y(X,F)
\end{tikzcd}}}
\diagram{\resizebox{\linewidth}{!}{\begin{tikzcd}[ampersand replacement=\&,column sep=1em]
\operatorname{Ext}^{r-p}_{\widehat{\underline O}_P}(\widehat Y;\widehat F,\widehat\Omega^r)
 \&\operatorname{Ext}^{r-p}(P,F,\Omega^r)\arrow[l]
 \&\varprojlim\operatorname{Ext}^{r-p}(P;F,J^{n}\Omega)\arrow[l]
 \&{}\arrow[l]
\end{tikzcd}}}
\[
\operatorname{Ext}^{r}
\]
% archive sheet number: 4
\scanpage{6}
\diagram{\resizebox{\linewidth}{!}{\fbox{\begin{tikzcd}[ampersand replacement=\&,column sep=1em,row sep=.35em]
\ldotp\ldotp\arrow[r]\&
 \operatorname{Ext}^{i}(X;F,G)\arrow[r]\&
 \operatorname{Ext}^{i}(X';F',G')\arrow[r]\&
 \operatorname{Ext}^{i+1}_{\unsure{-|a}}(X;F,G)\arrow[r]\&{}\\
\&\times\&\times\&\times\&\\
\ldots\&\operatorname{Ext}^{r-i}_{a}(X,G,F)\arrow[l]\&
 \operatorname{Ext}^{r-1-i}(X',\deleted{\unsure{F}}G',\deleted{\unsure{G}}F')\arrow[l]\&
 \operatorname{Ext}^{r-1-i}(X;G,F)\arrow[l]\&{}\arrow[l]
\end{tikzcd}}}}
\[
\begin{aligned}
\operatorname{Ext}^i_Y(X;F,G)&=\varinjlim\operatorname{Ext}^{i}(X;F_n,G)\\
\text{dual \unsure{à|de|a}}&\quad\varprojlim\operatorname{Ext}^{r-i}_a(X,G,F_n)
 \simeq\deleted{\unsure{\ell|E}}\operatorname{Ext}^{\unsure{r-i}}_{-}(\widehat X,\widehat G,\widehat F)
\end{aligned}
\]
\[
\begin{aligned}
\operatorname{Ext}^{\unsure{\bullet|i}}_{Y'}(X';F',G')&=\varinjlim\operatorname{Ext}^{\unsure{n|i}}(X';F'_n,G')\\
\text{dual à}&\quad\varprojlim\operatorname{Ext}^{r-1-i}(X',G',F'_n)\\
&\quad\operatorname{Ext}^{r-1-i}(\widehat X';\widehat G',\widehat F')\ ?\\
&\hspace{7em}\rotatebox{90}{$\simeq$}
\end{aligned}
\]
\usetikzlibrary{fit,shapes.geometric}
\diagram{\resizebox{\linewidth}{!}{\fbox{\begin{tikzcd}[ampersand replacement=\&,column sep=1.5em,row sep=.6em,
execute at end picture={\node[draw,ellipse,fit=(ymid)(ymidb)(ymidc),inner sep=0pt]{};
\node[draw,ellipse,fit=(yprime)(yprimeb),inner sep=1pt]{};}]
\ldotp\ldotp\arrow[r]\&\operatorname{Ext}^{i}_{-}(X,F,G)\arrow[r]\&
 |[alias=ymid]|{\operatorname{Ext}^i_Y(X;F,G)}\arrow[r]\&
 |[alias=yprime]|{\operatorname{Ext}^{i}_{Y'}(X';F',G')}\arrow[r,dash]\&\ldots\\
\&\&\otimes\&\otimes\&\\
-\&\operatorname{Ext}^{r-i}(\widehat X,\widehat G,\widehat F)\arrow[l]\&
 |[alias=ymidb]|{\operatorname{Ext}^{r-i}_a(\widehat X,\widehat G,\widehat F)}\arrow[l]\&
 |[alias=yprimeb]|{\operatorname{Ext}^{r-1-i}(\widehat X';\deleted{\unsure{F}}\widehat G',\widehat F')}\arrow[l]\&\ldotp\mathord{-}\arrow[l]\\
\&\rotatebox{90}{$\simeq$}\&\Vert\&\&\\
\&\operatorname{Ext}^{r-i}(X;G,F)\&
 |[alias=ymidc]|{\deleted{\operatorname{Ext}^{\unsure{r-i}}_{\unsure{a}}(X,G,F)}}\&\&
\end{tikzcd}}}}
\margin{plus général}

\diagram{\resizebox{\linewidth}{!}{\begin{tikzcd}[ampersand replacement=\&,column sep=1em]
\operatorname{Ext}^i_Y(X;F,G)\arrow[r]\&
 \operatorname{Ext}^{i}_{\deleted{\unsure{Y}}Z}(X;F,\overset{\unsure{\#|\ast}}{G})\&
 {\unsure{\mathord{-}|\to}}\&\operatorname{Ext}^{i}_{Z-Y}(X;F,G)
\end{tikzcd}}}
\diagram{\resizebox{.84\linewidth}{!}{\begin{tikzcd}[ampersand replacement=\&,column sep=1em]
\operatorname{Ext}^{r-1-i}_a(X_{/Y};G_{/Y},F_{/Y})\&
 \operatorname{Ext}^{r-1-i}_a(X_{/Z};G_{/Z},F_{/Z})\arrow[l]\&{}\arrow[l]
\end{tikzcd}}}
\margin{$\#$}
\[
\begin{aligned}
X_{/Y}&\longrightarrow X_{/Z}\\
\widehat{\vphantom X}&\hspace{4em}\widehat{\vphantom X}
\end{aligned}
\]
% archive sheet number: 5
\scanpage{7}
\mspage{\handwritten{- III -5-}}
\begingroup
\csname tagsleft@true\endcsname
le groupe des $p$-cochaînes \uline{alternées} de \leavevmode\rlap{C}\handwritten{\v{\phantom{C}}}ech, relativement au recouvrement $\underline U$, et à \deleted{\unsure{coefficients}} coefficients dans $\underline F$, il résulte de ce qui précède que l'on peut écrire
\[
C^p(\underline U,\underline F)=\mathop{\underset{\handwritten{\longrightarrow}}{\lim}}\limits_n C_n^p(M)\ .\tag*{(11)}
\]
Or, avec les notations du n\textsuperscript{o}1, $C_n^p(M)$ s'identifie à $K^{p+1}(\underline f,M)$, et l'application $\handwritten{\varphi}_{\unsure{nm|mn}}$ s'identifie à l'application $\handwritten{\varphi}_{\underline f^{n-m}}$ définie au n\textsuperscript{o}1. On a donc, pour tout $p\handwritten{\geq}0$, un isomorphisme canonique \deleted{\unsure{fonction}} fonctoriel en $M$
\[
C^p(\underline U,\underline F)\handwritten{\xrightarrow{\sim}}C^{p+1}((\underline f),M)\ .\tag*{(12)}
\]
La formule (3) du n\textsuperscript{o}1 montre en outre que ces isomorphismes sont compatibles avec les opérateurs cobord. Il en résulte ;

\uline{Proposition 2 .- Si} $X$ \uline{est} \inserted{\uline{un préschéma}} \uline{noethérien, ou} \deleted{\illegible{2}} \uline{un schéma don\unsure{t} l'espace de base est quasi-compact, il existe pour tout} $p\handwritten{\geq}1$ \uline{un isomorphisme canonique fonctoriel en} $\deleted{\unsure{M}}\handwritten{\underline F}$
\[
H^p(\underline U,\underline F)\handwritten{\xrightarrow{\sim}}H^{p+1}((\underline f),M)\ .\tag*{(13)}
\]
\uline{On a en outre une suite exacte fonctorielle en} $\deleted{\unsure{M}}\handwritten{\underline F}$
\[
0\handwritten{\to}H^0((\underline f),M)\handwritten{\to}M\handwritten{\to}H^0(\underline U,\underline F)\handwritten{\to}H^1((\underline f),M)\handwritten{\to}0\ .\tag*{(14)}
\]
Les relations (13) sont en effet évidentes d'après ce qui précède. D'autre part, \deleted{\unsure{on a} ${\unsure{H|C}}^{\unsure{0}}((\unsure{f}),\unsure{M})\unsure{=|,}{\unsure{H|C}}^{\unsure{1}}((\unsure{f}),\unsure{M})$} $M=C^0((\underline f),M)$ et \deleted{\illegible{1}} $C^0(\underline U,\underline F)={}=C^1((\underline f),M)$, donc $H^0(\underline U,\underline F)$ s'identifie au sous-groupe des 1-cocycles de $C^1((\underline f),M)$ ; la suite exacte (3) n'est autre alors que \inserted{celle qui résulte de} la définition des groupes de cohomologie $H^0((\underline f),M)$ et $H^1((\underline f),M)$.

\uline{Corollaire .- Supposons qu'il existe des} $g_i\handwritten{\in\Gamma}(U,O_X)$ \uline{tels que} $\handwritten{\sum_{i=1}^{\unsure{r|n}}}g_i f_{\unsure{i}}=1$ \inserted{\uline{au-dessus de} $U$}. \uline{Alors, pour tout faisceau quasi-cohérent} \deleted{\illegible{1}} $\deleted{\unsure{F}}\handwritten{\underline G}$ \uline{sur} $U$, \uline{on a} $H^i(\underline U,\underline G)=0$ \uline{pour} $i\handwritten{>}0$ ; \uline{si en outre} $X=U$, \uline{l'homomorphisme canonique} $M\handwritten{\to}H^0(\deleted{\unsure{M}}\underline{\unsure{U|M}},\underline F)$ \uline{est bijectif}. (les $X_{f_i}=U_i$ étant par hypothèse quasi-com\inserted{pacts)}

On peut en effet supposer $X=U$, et alors $H^i((\underline f),M)=0$ pour tout $i\handwritten{\geq\unsure{0|\alpha}}$ (n\textsuperscript{o}1, cor. de la prop.1); l'assertion découle alors aussitôt de (13) et (14). On observera que \inserted{puisque} $H^0(\underline U,\underline F)=H^{\deleted{\unsure{1}} \handwritten{0}}(U,\underline F)$, on redémontre ainsi le
\endgroup
% archive sheet number: 6
\scanpage{8}
\[
\begin{aligned}
&\deleted{\unsure{A|E|\operatorname{Ext}}}H^{r-i}(\widehat P,\Omega^{r}_{\deleted{X}P}(m))\\
&H^{\unsure{n|a}}(\widehat P,\Omega^{r}_{P}(m))\\
&\deleted{H^0(\widehat P\unsure{)|,)|,}\unsure{\mathord{-}|\mathord{=}}}\\
&\deleted{H^0(\widehat P,}
\end{aligned}
\]

\[
\resizebox{\linewidth}{!}{$\displaystyle\boxed{\begin{array}{cl}
\left.\left\{\begin{aligned}
&H^i(\widehat P,\Omega^r_{\widehat P}(m))=0 &&\text{si }0<i<n\\
&H^0(P,\Omega^r(m))\xrightarrow{\sim} H^0(\widehat P,\widehat\Omega^r(m))
\end{aligned}\right.\right|&
\Longleftrightarrow\ \begin{aligned}
&P-X\ \text{est de}\\[-2pt]
&\text{dim coh. }r-n\\[-2pt]
&[\text{et non plus}\\[-2pt]
&\quad\text{grand}]
\end{aligned}\\[3pt]
\begin{gathered}
\text{pour }m\text{ grand, ou tt }m\\
\text{(\unsure{au choix|en somme}) \unsure{on peut même}}\\[-2pt]
\text{remplacer }\Omega^r_P\text{ par }\underline O_P\ \unsure{\ldots|\text{ici}}
\end{gathered}&
\end{array}}$}
\]

Les schémas formels s'introduisent en Géométrie algébrique de la façon suivante :

1) Techniquement, pour énoncer et prouver de façon commode les deux th. fondamentaux des morphismes propres

% The source records the insertion at its destination. The box places its text
% below 3 as on the sheet; the arrow returns to the insertion point.
\begingroup
\newsavebox{\ScanEightItem}
\setbox\ScanEightItem=\vbox{\hsize=\linewidth\noindent
\tikz[remember picture,baseline=(sEightTwo.base)]{\node[draw,circle,inner sep=1pt] (sEightTwo) {2)};} \inserted{En théorie de Lefschetz - \unsure{Courant|Ce sont|Grauert}, \deleted{\unsure{pour|d'où}} \inserted{\unsure{le plus couramment des|le plus courant des}} \unsure{énoncés|énoncer} \deleted{les} \inserted{\unsure{donnant des|donnent des}} relations entre une variété (projective, ou locale) et son complété formel le long \unsure{``|d'}une section hyperplane''.}}

\noindent\tikz[remember picture,overlay]{\coordinate (sEightAt) at (.6,1.5ex);}3) En théorie formelle des Modules, pour \deleted{\unsure{pouvoir trouver|parfois trouver}} \inserted{\unsure{saisir}} une notion de familles modulaires\par
\smallskip\noindent\usebox{\ScanEightItem}\par
\noindent\tikz[remember picture,overlay]{\draw[->] (sEightTwo.west) .. controls +(-.6,.8) and +(-1.4,0) .. (sEightAt);}
\endgroup

4) En formalisme de dualité, les \unsure{duals|duels|énoncés} des groupes de cohomologie $H^p_Y(X,F)$ s'expriment par des $\operatorname{Ext}$ ou $H^{\unsure{*|q|p}}$ dans $\widehat X$ (complété formel le long de $Y$)

\margin{TSVP}
% archive sheet number: 7
\scanpage{9}
5) Théorie des groupes formels, \unsure{du|de|où} $\pi_1$ etc. \unsure{Ce|ce|la} sont \unsure{souvent} des schémas formels ....

De plus, notons

6) Dans la théorie des types de \uline{Greenberg}\unsure{\({}_{*}\)|.} - \uline{\unsure{Néron|Niron}} s'introduisent des schémas relatifs sur des espaces annelés qui ne sont pas nécessairement des préschémas ....
% archive sheet number: 8
